\section{The Meaning of Christmas}

At this time of year, many voices arise to tell us the ``true'' meaning of Christmas. On the one hand, we are told that it is a reminder to be ``nice'' to one another. Then there are the ``Marthas'' of the world who tell us to do less ``shopping'', apparently on the theory that it is immoral to want to give gifts to those we love.

For us, the question is whether there can be something absolutely \emph{new} born into the world, something higher that transcends the grim fatalism of space, time, matter, contingency. Whether the birth of a \textbf{New Man} is possible, likewise born of the \textbf{Holy Spirit} and \textbf{Holy Virgin}.

It is not a question of a new system, or a new religion, as though such abstractions have a life in themselves apart from the men and women who uphold them and are transformed by them. At Christmas, we are reminded of our real life, to be reintegrated in our consciousness, to seek to spiritually transform the world. And should we fail, we must try again.

With acknowledgment to a special person who pointed out Botticelli's painting \textit{Madonna of the Roses} to me and inspired this thought (wherever you are).


\flrightit{Posted on 2010-12-25 by Cologero}

\begin{center}* * *\end{center}

\begin{footnotesize}
\begin{sffamily}
\texttt{Matt on 2010-12-25 at 18:15 said:}

Cologero, I just discovered Meditations on the Tarot here yesterday and having been reading as much as I can. I will say this, my appreciation for the Christian doctrine is starting to increase quite a bit from reading it.

With that said, I hope you are having a very Merry Christmas.

\hfill

\texttt{Porfyreos on 2010-12-26 at 03:06 said:}

I also discovered Valentin Tomberg's ``Meditations on the Tarot'' recently, and it has changed my view on Christian doctrine in a favorable way as well. I am very grateful that this book is published here on this website's on-line library.

There seem to be great differences in spiritual approach and interpretation of history expressed in this book compared to the views expressed by Evola, Guénon and other ``traditionalist'' writers. Valentin Tomberg has a much more favorable attitude toward thinkers like Bergson, Jung, Teilhard de Chardin, as well as Steiner's anthroposophy, Blavatsky's theosophy, and the reality and importance of reincarnation and evolution. It might be very fruitful to contemplate these areas of disagreements in order to purify Tomberg's Christian Hermetic path, or, if this path cannot hold the ground against traditionalist criticism, to save what can be saved of Tomberg's insights and bring them over to a better starting point.

Until I read ``Meditations on the Tarot``, Evola's ``magical hero'' was my initiatic ideal, now I see the ideal of the ``religious saint'' much clearer, and I wonder which of these spiritual ideals should be the guiding star. May be a hermetic synthesis between these two tendencies is needed in order to create ``A New Man''.

I wish a merry Christmas to all the contributors on this thought-provoking website.

\hfill

\texttt{Cologero on 2010-12-26 at 12:04 said:}

``Meditations on the Tarot'' is a work of spirituality, not metaphysics, so it deals with the phenomenal world rather than principles. As to the writers you mention, notice what Tomberg writes about in the Moon. It does not imply unreserved acceptance. Gnosis requires a change in one's state of being. So even if a thinker is not completely orthodox (in the metaphysical senses), his thought can be hermetically transformed to effect such a change, just as a base metal is transformed into gold. So, it is more a question of perspective rather than a contradiction. Evola is not so pure as Guenon; there are many non-Traditional thinkers that Evola refers to and build on.

The ``magical hero'' may be the appropriate model for some, depending on individual characteristics and interests; after all, the purification of the Will is part of the Hermetic process. Tomberg meditates on the Tarot and recommends meditating on the symbolism of the Gospels. However, the same method could apply to meditations on other Traditional symbols. The goal is depth, not presenting the correct appearances. Fanaticism in religious beliefs is incompatible with esoterism.

\hfill

\texttt{Graham on 2010-12-28 at 19:31 said:}

Porfyreos, Matt-

Have you read the Philokalia? It might surprise you.

\hfill

\texttt{Matt on 2010-12-28 at 20:12 said:}

Graham,

The library at the university I attend has all of the Philokalia. So far I have read St. Anthony's contribution. My plan is to read at least 2 sections a week once school starts up again.


\hfill

\end{sffamily}
\end{footnotesize}

